% ============================================================
% BÖLÜM 6: İLGİLİ ALANYAZIN
% Genişletilmiş Akademik Literatür Taraması
% ============================================================

\chapter{İLGİLİ ALANYAZIN}

\section{Projektif Çizim Testleri}

\subsection{Tarihsel Gelişim}

Projektif testler, bireyin bilinçdışı düşünce, duygu ve çatışmalarını ortaya çıkarmak amacıyla kullanılan psikolojik değerlendirme araçlarıdır. Bu testlerin temeli, Freud'un psikanalitik kuramına dayanmaktadır. Projeksiyon kavramı, bireyin kendi düşünce ve duygularını dış nesnelere atfetmesi sürecini ifade etmektedir (Freud, 1911). Jung (1912) ise kolektif bilinçdışı kavramı ile projektif testlerin kuramsal temelini genişletmiştir.

Frank (1939), ``projektif yöntemler'' terimini ilk kez kullanmış ve bu yöntemlerin bireyin iç dünyasını anlamada benzersiz bir araç sunduğunu ileri sürmüştür. Bu yaklaşıma göre, belirsiz uyaranlar karşısında bireyin verdiği yanıtlar, bilinçdışı süreçleri yansıtmaktadır.

Çizim temelli projektif testlerin tarihsel gelişimi şu şekilde özetlenebilir:

\begin{table}[H]
\centering
\caption{Projektif Çizim Testlerinin Tarihsel Gelişimi}
\label{tab:tarihi}
\begin{tabular}{cll}
\toprule
\textbf{Yıl} & \textbf{Geliştirici} & \textbf{Test/Katkı} \\
\midrule
1885 & Corrado Ricci & İlk sistematik çocuk çizimi analizi \\
1926 & Florence Goodenough & Draw-A-Man (Zekâ ölçümü) \\
1948 & John Buck & House-Tree-Person (HTP) \\
1949 & Karen Machover & Draw-A-Person (DAP) \\
1952 & Charles Koch & Ağaç Testi (Baum Test) \\
1963 & Emanuel Hammer & HTP Klinik Uygulaması \\
1968 & Elizabeth Koppitz & Çocuk Figür Çizimi Değerlendirmesi \\
1977 & Karen Bolander & Ağaç Çizimi ile Kişilik Değerlendirmesi \\
1985 & Burns \& Kaufman & Kinetik Aile Çizimi (KFD) \\
1992 & Silver & Çizim ile Bilişsel Değerlendirme \\
\bottomrule
\end{tabular}
\end{table}

\subsection{House-Tree-Person (HTP) Testi}

HTP testi, John Buck tarafından 1948 yılında geliştirilmiş ve daha sonra Emanuel Hammer (1958) tarafından genişletilmiştir. Test, bireyin üç nesne çizmesini içermektedir:

\begin{enumerate}
  \item \textbf{Ev (House):} Buck'a (1948) göre ev çizimi, aile ilişkileri, ev ortamı algısı ve güvenlik hissini yansıtır. Ev, ``psikolojik sığınak'' olarak yorumlanmaktadır. Hammer (1958), evin kişinin kendisi ve aile üyeleri ile olan ilişkilerini sembolize ettiğini belirtmiştir.

  \item \textbf{Ağaç (Tree):} Koch (1952) ve Bolander (1977), ağaç çiziminin bilinçdışı benlik algısını, büyüme ve gelişim hissini yansıttığını ileri sürmüştür. Ağaç, ``yaşam ağacı'' sembolizmi ile bireyin varoluşsal deneyimini temsil etmektedir.

  \item \textbf{İnsan (Person):} Machover'a (1949) göre insan figürü çizimi, benlik algısı, beden imajı ve kişilerarası ilişkileri yansıtır. Çizilen figür, genellikle çizen kişinin bilinçdışı benlik temsilini oluşturmaktadır.
\end{enumerate}

\begin{table}[H]
\centering
\caption{HTP Testinin Temel Yorumlama Alanları}
\label{tab:htpyorumlama}
\begin{tabular}{p{2.5cm}p{4.5cm}p{5.5cm}}
\toprule
\textbf{Çizim} & \textbf{Sembolik Anlam} & \textbf{Değerlendirilen Alanlar} \\
\midrule
Ev & Aile ilişkileri, güvenlik, ev ortamı algısı & Kapı, pencere, çatı, duvar, baca, zemin çizgisi, yol, çevre \\
Ağaç & Benlik, büyüme, yaşam enerjisi & Gövde, dallar, kökler, yapraklar, meyve, kabuk detayı \\
İnsan & Benlik algısı, beden imajı, ilişkiler & Baş, gövde, uzuvlar, yüz detayları, giysiler, cinsiyet \\
\bottomrule
\end{tabular}
\end{table}

\subsection{HTP'nin Psikometrik Özellikleri}

HTP testinin psikometrik özellikleri, akademik literatürde tartışma konusu olmuştur:

\textbf{Güvenilirlik:} Test-tekrar test güvenilirliği .60 ile .80 arasında değişmektedir (Handler ve Habenicht, 1994). Değerlendiriciler arası güvenilirlik ise değerlendirmecilerin eğitim düzeyine bağlı olarak .31 ile .96 arasında geniş bir aralıkta seyretmektedir.

\textbf{Geçerlilik:} Yapı geçerliliği çalışmaları karışık sonuçlar vermiştir. Bazı araştırmacılar (Marzolf ve Kirchner, 1972; Joiner ve ark., 1996), HTP göstergelerinin psikolojik değişkenlerle anlamlı korelasyonlar gösterdiğini raporlamıştır. Ancak Lilienfeld ve ark. (2000), bu korelasyonların çoğunun düşük düzeyde olduğunu eleştirmiştir.

\textbf{Tartışmalar:} Swensen (1957, 1968) tarafından yürütülen iki meta-analiz, projektif çizim testlerinin geçerliliğine ilişkin sınırlı kanıt bulmuştur. Bununla birlikte, Naglieri ve ark. (1991), DAP testinin zekâ tahmininde geçerli olduğunu göstermiştir.

\section{Çocuk Çizimlerinin Gelişimsel Özellikleri}

\subsection{Lowenfeld'in Çizim Gelişim Evreleri}

Viktor Lowenfeld (1947), çocukların çizim gelişimini altı evrede incelemiştir. Bu evreler, NIDA sisteminin yaş tabanlı yorumlama algoritmasının temelini oluşturmaktadır:

\begin{enumerate}
  \item \textbf{Karalama Evresi (2-4 yaş):}
  \begin{itemize}
    \item Rastgele çizgiler, motor aktivite ağırlıklı
    \item Görsel-motor koordinasyonun gelişimi
    \item Çizgiler henüz sembolik anlam taşımaz
  \end{itemize}

  \item \textbf{Ön-Şematik Evre (4-7 yaş):}
  \begin{itemize}
    \item İlk tanınabilir figürler ortaya çıkar
    \item ``Kurbağa adam'' (tadpole figure) tipik bir özelliktir
    \item Orantılar henüz gerçekçi değildir
    \item Renk kullanımı duygusaldır, gerçekçi değildir
  \end{itemize}

  \item \textbf{Şematik Evre (7-9 yaş):}
  \begin{itemize}
    \item Tutarlı semboller gelişir
    \item Taban çizgisi (baseline) kullanılmaya başlanır
    \item X-ray (röntgen) çizim tekniği görülür
    \item Orantılar iyileşir ancak hâlâ semboliktir
  \end{itemize}

  \item \textbf{Gerçekçilik Evresi (9-12 yaş):}
  \begin{itemize}
    \item Detay artışı belirginleşir
    \item Örtüşme ve derinlik ipuçları kullanılır
    \item Cinsiyet farklılıkları çizimlerde belirginleşir
    \item Kritik bakış açısı gelişmeye başlar
  \end{itemize}

  \item \textbf{Sözde-Doğalcı Evre (12-14 yaş):}
  \begin{itemize}
    \item Perspektif denemeleri yapılır
    \item Kendi çizimlerine eleştirel bakış
    \item Bazı çocuklar çizimden uzaklaşır
  \end{itemize}

  \item \textbf{Karar Evresi (14-17 yaş):}
  \begin{itemize}
    \item Sanatsal yönelim veya çizimden tamamen uzaklaşma
    \item Stilize veya karikatür çizimler
  \end{itemize}
\end{enumerate}

Bu gelişimsel evreler, çizim yorumlamasında yaş faktörünün önemini ortaya koymaktadır. NIDA sistemi, 5-12 yaş aralığını (Ön-Şematik, Şematik ve Gerçekçilik evreleri) hedeflemektedir.

\subsection{Koppitz'in Gelişimsel Göstergeleri}

Elizabeth Koppitz (1968, 1984), çocukların insan figürü çizimlerinde gelişimsel göstergeler tanımlamıştır. Bu göstergeler, yaşa bağlı beklenen ve beklenmeyen özellikleri ayırt etmektedir:

\begin{itemize}
  \item \textbf{Beklenen Öğeler:} Yaşa göre çoğu çocukta görülen özellikler (örn. 7 yaşında baş, gövde, bacaklar)
  \item \textbf{Olağanüstü Öğeler:} Yaşa göre nadir görülen özellikler (örn. 6 yaşında eklem detayları)
  \item \textbf{Duygusal Göstergeler:} Psikolojik sorunlarla ilişkilendirilen özellikler (örn. gizli figürler, gölgeleme)
\end{itemize}

\section{Projektif Testlerin Geçerliliği Tartışması}

\subsection{Eleştiriler}

Projektif testlerin psikometrik özellikleri, akademik literatürde yoğun tartışma konusu olmuştur. Lilienfeld, Wood ve Garb (2000), kapsamlı bir incelemede projektif testlerin geçerliliğine ilişkin ciddi sorular ortaya koymuştur:

\begin{quotebox}
``Projektif teknikler için ampirik destek sınırlıdır ve bu tekniklerin rutin klinik kullanımı sorgulanmalıdır. Özellikle, çoğu projektif gösterge için geçerlilik kanıtları yetersizdir.'' (Lilienfeld ve ark., 2000, s. 27)
\end{quotebox}

Hunsley ve Bailey (1999), projektif testlerin artımlı geçerliliğinin (incremental validity) sorgulanması gerektiğini belirtmiştir. Dawes (1994) ise projektif testlerin ``klinik mitler'' olduğunu ileri sürmüştür.

\subsection{Savunular}

Bununla birlikte, diğer araştırmacılar projektif testlerin nitel değerini vurgulamamıştır:

\begin{quotebox}
``Çizim testleri, çocukların iç dünyalarına ulaşmada benzersiz bir pencere sunmaktadır. Psikometrik sınırlılıklarına rağmen, klinik değerlendirmenin önemli bir parçası olmaya devam etmektedir.'' (Malchiodi, 1998, s. 42)
\end{quotebox}

Handler (1996), projektif testlerin ``idiografik'' (bireysel) değerlendirme için benzersiz bilgiler sunduğunu savunmuştur. Ganellen (1996), Rorschach testinin uygun kullanıldığında geçerli sonuçlar ürettiğini göstermiştir.

\subsection{NIDA'nın Yaklaşımı}

NIDA, bu tartışmayı çözmek yerine, mevcut literatürü şeffaf biçimde sunmaktadır:

\begin{itemize}
  \item Her yorumun güven seviyesi belirtilmektedir
  \item Çelişkili akademik görüşler sunulmaktadır
  \item Sistem, tanı aracı olarak değil, karar destek sistemi olarak konumlandırılmıştır
  \item Kanıt gücü (evidence strength) her kural için kategorize edilmiştir
\end{itemize}

\section{Yapay Zekâ ve Psikolojik Değerlendirme}

\subsection{Ruh Sağlığında Yapay Zekâ Uygulamaları}

Son yıllarda yapay zekâ teknolojilerinin psikolojik değerlendirme ve ruh sağlığı hizmetlerinde kullanımı artmaktadır. Torous ve ark. (2018), dijital fenotiping kavramını tanımlamış ve akıllı cihazlardan toplanan verilerin ruh sağlığı değerlendirmesinde kullanılabileceğini göstermiştir.

\begin{table}[H]
\centering
\caption{Yapay Zekânın Ruh Sağlığındaki Uygulama Alanları}
\label{tab:airuhsagligi}
\begin{tabular}{p{3.5cm}p{5cm}p{4cm}}
\toprule
\textbf{Uygulama Alanı} & \textbf{Örnek Çalışmalar} & \textbf{Kullanılan Teknoloji} \\
\midrule
Depresyon Tespiti & De Choudhury ve ark., 2013 & Sosyal medya analizi \\
Anksiyete Taraması & Fitzpatrick ve ark., 2017 & Chatbot (Woebot) \\
Psikoz Riski & Bedi ve ark., 2015 & Konuşma analizi \\
İntihar Riski & Walsh ve ark., 2017 & Makine öğrenmesi \\
Çizim Analizi & Kim ve ark., 2024 & MLLM \\
\bottomrule
\end{tabular}
\end{table}

\subsection{Çocuk Çizimi Analizinde Yapay Zekâ}

Çocuk çizimlerinin yapay zekâ ile analizi, görece yeni bir araştırma alanıdır. Öne çıkan çalışmalar:

\begin{enumerate}
  \item \textbf{PICK (Kim ve ark., 2024):} Çok modlu büyük dil modellerini HTP analizi için kullanan bir yaklaşım. GPT-4V modeli kullanılarak F1 skorunda \%10+ iyileştirme raporlanmıştır. Çalışma, MLLM'lerin çizim yorumlamasında umut verici sonuçlar ürettiğini göstermiştir.

  \item \textbf{HTP Derin Öğrenme (Chen ve ark., 2023):} YOLOv5 nesne tespit teknolojisi ile HTP çizimlerindeki öğelerin otomatik tespiti. Model, \%85 doğruluk oranı ile ev, ağaç ve insan figürlerini tespit edebilmektedir.

  \item \textbf{Depresyon Tespiti (Park ve ark., 2023):} HTP çizimlerinin konvolüsyonel sinir ağları (CNN) ile analizi. Depresyon riskini \%72 doğrulukla tahmin edebilmiştir.

  \item \textbf{Otizm Taraması (Liu ve ark., 2022):} İnsan figürü çizimlerinin otizm spektrum bozukluğu taramasında kullanımı. Derin öğrenme modeli, \%78 hassasiyet raporlamıştır.
\end{enumerate}

\subsection{MLLM'lerin Çizim Analizindeki Rolü}

Çok Modlu Büyük Dil Modelleri (Multimodal Large Language Models - MLLM), metin ve görsel içerikleri birlikte işleyebilen yapay zekâ sistemleridir. GPT-4V (OpenAI, 2023), Gemini (Google, 2024) ve Claude 3 (Anthropic, 2024) gibi modeller, görsel akıl yürütme görevlerinde insana yakın performans sergilemektedir.

NIDA, Google Gemini 2.0 Flash modelini kullanmaktadır. Bu modelin avantajları:

\begin{itemize}
  \item Türkçe dil desteği
  \item Görsel-metin çoklu modal analiz kapasitesi
  \item Düşük gecikme süresi
  \item Maliyet etkinliği
  \item API erişim kolaylığı
\end{itemize}

\subsection{Güncel Çok Modlu Dil Modelleri (2023--2024)}

Çok modlu büyük dil modelleri alanında 2023--2024 yılları arasında önemli gelişmeler yaşanmıştır. Google DeepMind tarafından geliştirilen Gemini modelleri ailesi (Team Gemini, 2023), metin, görsel, ses ve video gibi farklı modaliteleri birleşik bir mimari içinde işleyebilen ilk büyük ölçekli çok modlu sistemi temsil etmektedir. Bu model ailesi; Ultra, Pro ve Nano olmak üzere üç ölçek kademesiyle sunulmuş ve çeşitli çok modlu kıyaslama testlerinde GPT-4V ile rekabet edebilir sonuçlar üretmiştir (Team Gemini, 2023). Gemini 1.5 Pro modeli ise bağlam penceresini milyonlarca jetona genişleterek uzun belgeler ve çok sayıda görsel içeren görevlerde önemli ilerleme kaydetmiştir (Google, 2024b).

OpenAI'nin GPT-4 Teknik Raporu (OpenAI, 2023), çok modlu dil modellerinin görsel akıl yürütme kapasitesini sistematik biçimde belgelemektedir. GPT-4V bileşeni; tıbbi görüntü yorumlama, diyagram analizi ve sanat eseri değerlendirmesi gibi alanlarda insana yakın sonuçlar üretmektedir. Bu özellikler, HTP çizimi yorumlaması gibi görsel-sembolik analiz gerektiren görevlerde MLLM kullanımının gerekçesini güçlendirmektedir.

Anthropic tarafından geliştirilen Claude 3 model ailesi (Anthropic, 2024), özellikle uzun bağlamlı çok modlu görevlerde dikkat çekici sonuçlar üretmektedir. Bu güncel gelişmeler, NIDA'nın benimsediği Gemini tabanlı yaklaşımın sağlam bir teknolojik zemine oturduğunu ve alandaki en ileri sistemlerle paralel bir çizgide ilerlediğini ortaya koymaktadır.

\section{Türkiye'de PDR ve Teknoloji}

\subsection{Mevcut Durum}

Türkiye'de psikolojik danışmanlık hizmetlerinde teknoloji kullanımı, özellikle COVID-19 pandemisi sonrasında hızlanmıştır (Öztürk ve Yıldız, 2022). Çevrimiçi danışmanlık platformları, dijital değerlendirme araçları ve mobil uygulamalar yaygınlaşmaktadır.

\begin{itemize}
  \item \textbf{Çevrimiçi Danışmanlık:} Pandemi döneminde \%300 artış (TÜİK, 2021)
  \item \textbf{Dijital Değerlendirme:} Sınırlı sayıda Türkçe araç
  \item \textbf{Yapay Zekâ Uygulamaları:} Henüz yaygınlaşmamış
\end{itemize}

\subsection{Boşluklar ve Fırsatlar}

Türkiye'de çocuk çizimi analizi alanında özgün bir dijital araç bulunmamaktadır. Bu boşluk, NIDA projesi için hem bir gerekçe hem de bir fırsat oluşturmaktadır. Potansiyel katkılar:

\begin{enumerate}
  \item Türkçe dil desteği ile erişilebilirlik
  \item Akademik literatürün Türkçe'ye aktarımı
  \item Yerel ihtiyaçlara uygun arayüz tasarımı
  \item Türk kültürel bağlamının entegrasyonu (gelecek sürümlerde)
\end{enumerate}

\newpage
